% Projeto de pesquisa ABNT (abnTeX2, NBR 15287). Marcadores <<CHAVE>> sao trocados por doclib.fill_template().
% Sem resumo; secoes padrao de projeto. Compile com doclib.latex_abnt("projeto.tex").
\documentclass[12pt, oneside, a4paper, chapter=TITLE, section=TITLE, sumario=abnt-6027-2012, english, brazil]{abntex2}

\usepackage{lmodern}
\usepackage[T1]{fontenc}
\usepackage[utf8]{inputenc}
\usepackage{indentfirst}
\usepackage{color}
\usepackage{graphicx}
\usepackage{microtype}
\usepackage{booktabs}
\usepackage[alf]{abntex2cite}
% Links sem caixas coloridas (impressao e PDF entregue).
\hypersetup{hidelinks}

\titulo{<<TITULO>>}
\autor{<<AUTOR>>}
\local{<<LOCAL>>}
\data{<<ANO>>}
\orientador{<<ORIENTADOR>>}
\instituicao{<<INSTITUICAO>>}
\tipotrabalho{Projeto de pesquisa}
\preambulo{<<PREAMBULO>>}

\setlength{\parindent}{1.25cm}
\setlength{\parskip}{0pt}
\SingleSpacing
\OnehalfSpacing
\makeindex

\begin{document}
\selectlanguage{brazil}
\frenchspacing

\imprimircapa
\imprimirfolhaderosto*

\pdfbookmark[0]{\contentsname}{toc}
\tableofcontents*
\cleardoublepage

\textual

\chapter{Introdução}
<<INTRODUCAO>>

\section{Problema de pesquisa}
<<PROBLEMA>>

\section{Objetivos}
\subsection{Objetivo geral}
<<OBJETIVO_GERAL>>
\subsection{Objetivos específicos}
<<OBJETIVOS_ESPECIFICOS>>

\section{Justificativa}
<<JUSTIFICATIVA>>

\chapter{Referencial teórico}
<<REFERENCIAL>>

\chapter{Metodologia}
<<METODOLOGIA>>

\chapter{Cronograma}
\begin{table}[htb]
\ABNTEXfontereduzida
\caption{Cronograma de atividades}
\centering
\begin{tabular}{lcccc}
\toprule
Atividade & Mês 1 & Mês 2 & Mês 3 & Mês 4 \\
\midrule
<<CRONOGRAMA_LINHAS>>
\bottomrule
\end{tabular}
\fonte{Elaborado pelo autor (<<ANO>>).}
\end{table}

\postextual
\bibliography{referencias}

\end{document}
